\documentclass[11pt]{article}
\usepackage[utf8]{inputenc}
\usepackage[T1]{fontenc}
\usepackage{amsmath,amssymb,amsthm}
\usepackage{booktabs}
\usepackage{tabularx}
\usepackage{graphicx}
\usepackage[margin=2.5cm]{geometry}
\usepackage[round]{natbib}
\usepackage{hyperref}
\hypersetup{colorlinks=true, linkcolor=blue, citecolor=blue, urlcolor=cyan}

\newtheorem{theorem}{Theorem}
\newtheorem{corollary}{Corollary}
\newcommand{\inn}[2]{\langle #1, #2 \rangle}

\title{\bfseries The Projection Law: Distilling 40 Years of Molecular Dynamics into a Universal Correction Operator}
\author{Alex Welcing\\ \small Lupine Science, Union City, NJ\\ \small \texttt{alexwelcing@gmail.com}\\ \small ORCID: \href{https://orcid.org/0009-0002-1602-8545}{0009-0002-1602-8545}}
\date{June 29, 2026}

\begin{document}
\maketitle

\begin{abstract}
\noindent
When many independently constructed models agree, the agreement is routinely
read as confidence. We formalize and test the opposite reading: a model family
is a projection operator, fitting drives every member toward the point of the
family's reachable set nearest the truth, and the shared residual --- one
direction in observable space --- is a fingerprint of whatever constraint binds
the family. We prove the core as machine-checked theorems in Lean~4 (seven
theorems, zero \texttt{sorry}): best approximations are unique and share one
residual lying in the family's normal cone; the participation ratio of the
error second moment is a closed-form gauge of the systematic fraction; and the
empirical second moment concentrates entrywise. We test the law's sharpest
consequence --- errors organize by constraint, not by implementation --- at
three layers of one epistemic stack: classical interatomic potentials (559
models), foundation MLIPs ($4{\times}2$ MatPES factorial), and DFT
implementations (12 ACWF methods). We then report a new Round-2 3$\times$3$\times$3
elastic-constant benchmark of 16 cubic metals with four MatPES foundation MLIPs.
A one-vector-per-functional correction operator, validated with leave-one-out
cross-validation, reduces the benchmark mean absolute error from $17.84$~GPa to
$10.36$~GPa with zero no-harm violations, improving every model on both PBE and
approximate r$^2$SCAN targets. The formal pre-registered hypothesis tests on
this benchmark show that the functional-clustering predictions do not survive
in the 3d/4d subset, while the operator-vs-ensemble head-to-head passes on MAE
but not on conformal coverage. These mixed results define the next experimental
step: a class-aware operator and a direct A6 bridge test, not a larger
single-bias claim.
\end{abstract}

\section{The HPC Bottleneck: Ensembles and DFT Correction Loops}
\label{sec:hpc}

Molecular dynamics (MD) simulations at operational scale --- materials
discovery pipelines, alloy design campaigns, battery electrolyte screening ---
routinely consume $10^6$--$10^9$ core-hours per study. Two patterns dominate
the cost structure.

\textbf{Ensemble averaging.} To estimate uncertainty, practitioners run
multiple interatomic potentials or machine-learned interatomic potentials
(MLIPs) and pool their predictions. The practice is standard in climate
modeling \citep{pirtle2010,parker2011,bishop2013}, machine learning
\citep{fort2019,lakshminarayanan2017}, and materials simulation
\citep{frederiksen2004,lejaeghere2016}. Its cost is linear in the ensemble
size: $N$ models, $N$ times the compute. The implicit assumption is that
agreement among models signals reliability.

\textbf{DFT correction loops.} To recover accuracy, practitioners run
expensive density-functional-theory (DFT) calculations on a subset of
configurations and correct the cheaper potential's predictions iteratively.
The cost is dominated by the DFT calls; the correction is typically a scalar
shift or a linear fit, with no systematic transfer to new models or new
elements.

Both bottlenecks share a common root: the treatment of model error as
isotropic noise. If errors were isotropic, ensemble averaging would reduce
variance and DFT correction would be a local regression. But if errors are
\emph{structured} --- if they share a low-dimensional bias direction determined
by the modeling constraint --- then ensemble averaging leaves the bias
intact, and DFT correction can be distilled into a single, transferable
operator.

This paper establishes that structured-error regime as a theorem, verifies it
across three layers of one epistemic stack, and operationalizes it as a
software operator that replaces both bottlenecks with a single correction
call.

\section{Theory: Projection Operators and the Normal Cone}
\label{sec:theory}

The mathematics of best approximation is classical
\citep{deutsch2001,kinderlehrer1980}; we claim only its epistemic application
to model ensembles and the machine-checked instantiation chain. Throughout,
$E$ is a real inner-product space (observable space), $s \subseteq E$ a model
family's reachable set, $T \in E$ the truth.

\begin{theorem}[Normal-cone criterion]\label{thm:cone}
Let $s$ be convex and $p$ a best approximation of $T$ in $s$
(i.e.\ $p \in s$ and $\|T-p\| \le \|T-q\|$ for all $q \in s$). Then the
residual $T - p$ lies in the normal cone of $s$ at $p$:
$\inn{T-p}{q-p} \le 0$ for every $q \in s$.
\end{theorem}

\begin{theorem}[Consensus]\label{thm:consensus}
Best approximations onto a convex family are unique; consequently any two best
approximations of the same target share an \emph{identical} residual. The
residual --- hence any agreement among independently fitted members --- is
determined by the (family, target) pair alone, and vanishes exactly when the
truth lies in the family.
\end{theorem}

Agreement among models sharing a constraint is therefore a measurement of the
constraint, not evidence of truth: the formal content of the qualitative
warnings of \citet{parker2011} and \citet{pirtle2010}, and a strengthening of
error-correlation dependence measures \citep{bishop2013} in that, combined
with the factorial designs of \S\ref{sec:mlip}--\ref{sec:dft}, it identifies
\emph{which} constraint binds.

\begin{theorem}[Gauge]\label{thm:gauge}
Model errors $e = b + \xi$ with shared bias $b$ and isotropic noise of scale
$\sigma$ in $d$ dimensions have error second moment with one eigenvalue
$\|b\|^2 + \sigma^2$ (bias axis) and $d{-}1$ eigenvalues $\sigma^2$, and
participation ratio
\[
\mathrm{PR}(d,\rho) \;=\; \frac{(\rho+d)^2}{(\rho+1)^2 + (d-1)},
\qquad \rho = \|b\|^2/\sigma^2 ,
\]
with $\mathrm{PR}(d,0)=d$, $1 \le \mathrm{PR} \le d$, strict decrease in
$\rho$, and the explicit collapse bound
$\mathrm{PR}-1 \le 3(d-1)/\rho$ for $\rho \ge d$.
\end{theorem}

The algebra is a one-line corollary of spiked-covariance theory
\citep{johnstone2001} and the participation ratio is a standard
effective-dimension metric \citep{gao2015}; the value here is the consilience
it enables (\S\ref{sec:zerok}) and the machine-checked monotonicity that
licenses inverting measured PR into a systematic fraction
$\alpha = \rho/(\rho+1)$.

\begin{theorem}[Ribbon/consensus decoupling]\label{thm:decouple}
The participation ratio of a shared-axis ensemble $\{c_i u\}$ is $1$ for every
sign pattern of the coefficients $c_i$, while mean pairwise alignment
distinguishes them (e.g.\ $1$ vs.\ $-1/3$ for three members). PR detects the
\emph{axis}; alignment detects \emph{sign coherence}; they are distinct order
parameters.
\end{theorem}

\begin{theorem}[Affine decomposition]\label{thm:affine}
Let $K = a + L$ be a closed affine reachable set and $p^*$ the best
approximation of $T$ in $K$. For any fitted $p \in K$, the residual decomposes
as
\[
T - p \;=\; b \, +\, \xi(p),
\qquad b = T - p^* \in L^\perp, \quad \xi(p) = p^* - p \in L,
\]
with $b \perp \xi(p)$.
\end{theorem}

Theorem~\ref{thm:affine} turns the bias-plus-within-family \emph{split} of
Theorem~\ref{thm:gauge} from a modeling assumption into a derivable
consequence for affine families.

\begin{theorem}[Local normal cone, smooth non-convex families]\label{thm:smooth}
Let $f: \mathbb{R}^k \to E$ be a $C^1$ immersion and $x^*$ a local minimizer
of $\|T - f(x)\|$. Then the residual $T - f(x^*)$ is orthogonal to the tangent
space $\operatorname{range} Df(x^*)$.
\end{theorem}

Theorem~\ref{thm:smooth} licenses testing tangent-space orthogonality for
neural-network reachable sets: the conclusion is pointwise (one normal space
per local minimizer), not the global consensus theorem. The empirical
clustering in \S\ref{sec:mlip} is interpreted as nearby local minimizers
landing in a common normal direction; it does not follow from
Theorem~\ref{thm:smooth} that every local minimizer shares the same residual.

\begin{theorem}[Finite-sample concentration]\label{thm:concentration}
Let $X_1, \dots, X_n$ be i.i.d.\ bounded random vectors in $\mathbb{R}^d$
($\|X_k\|_\infty \le B$) and let $M$ and $\widehat M_n$ be the population and
empirical entrywise second-moment matrices. For every entry $(i,j)$ and
$\varepsilon > 0$,
\[
\mathbb{P}\bigl(|\widehat M_{n,ij} - M_{ij}| \ge \varepsilon\bigr)
\;\le\; 2 \exp\!\bigl(-n\varepsilon^2 / (2B^4)\bigr).
\]
Moreover, the participation ratio
$\mathrm{PR}(M) = (\operatorname{tr} M)^2 / \operatorname{tr}(M^2)$ is
continuous wherever the denominator is non-zero.
\end{theorem}

All seven theorems are verified in Lean~4 against a pinned Mathlib as part of
the program's larger formal corpus --- 225 theorem and lemma declarations
across the current build, checked by \texttt{lake build} with zero
\texttt{sorry} and zero new axioms --- including a conditional hyper-ribbon
bound (spectral-decay hypotheses $\Rightarrow \mathrm{PR}<2$), a
context-correction theorem bundle whose ``correction closes the residual''
statement is the formal counterpart of the out-of-sample calibration result
below, and a formal audit module that previously caught an overstatement in
the program's own empirical work.

\section{0~K Verification: LAMMPS PR Data and Bias Vectors}
\label{sec:zerok}

The observational base \citep{welcing2026instance}: elastic-constant
($C_{11}, C_{12}, C_{44}$) errors of 559 classical potentials across 15 cubic
metals, computed at 0~K with a single LAMMPS harness. Three facts matter for
the law.

\textbf{Low, stationary error dimensionality.} Participation ratios of the
42 multi-element potentials occupy $[1.00, 2.29]$ out of 3 with median $1.09$
(dataset pinned in the replication kit), with no trend across forty years of
potential development --- accuracy improved; geometry did not. You cannot fit
your way off a constraint surface.

\textbf{Within-family consensus.} Within a functional-form family the errors are
nearly identical (within-family reference--prediction correlation $r = 0.95$
across families vs.\ $0.82$ pooled). The shared residual is a property of the
family, not of any individual parameterization.

\textbf{Gauge consistency.} The median PR of $1.09$ inverts
(Theorem~\ref{thm:gauge}, $d{=}3$) to a systematic fraction $\alpha \approx 0.98$,
in agreement with the within-family correlation ($0.95$) and the rank-one share
($0.96$). These three diagnostics are functionals of closely related
second-moment structure on overlapping data --- under the bias-plus-noise model
they are algebraically coupled --- so their agreement is internal consistency,
not independent corroboration.

The binding constraint at this layer is the functional form itself; the oldest
example is exact: central-force pair potentials satisfy the Cauchy relation
$C_{12}=C_{44}$ identically \citep{born1954}, so for any material violating
it, every pair potential's error necessarily contains the same forced
component --- the constraint dictates the direction.

We emphasize the distinction from parameter-space sloppiness: hyper-ribbon
geometry was introduced for the parameter manifolds of \emph{single} models
\citep{transtrum2011,machta2013,quinn2023,kurniawan2022}. The object here is
different --- error vectors of an ensemble of \emph{distinct} models in
observable space --- though the two pictures are related through the
projection of parameter manifolds into data space.

\section{The Operator: \texttt{lupine.correct()} + Conformal UQ}
\label{sec:operator}

The theoretical and empirical results above license a practical operator. The
shared error is one direction in observable space; one correction vector per
(element, observable) repairs every model in the family at once. The operator
is implemented as \texttt{lupine.correct()} with the following contract:

\textbf{Input.} A model prediction $y \in \mathbb{R}^d$ for a target
observable (e.g.\ elastic constants $C_{11}, C_{12}, C_{44}$); an element
identifier; a family tag (e.g.\ \texttt{classical}, \texttt{mlip-pbe},
\texttt{mlip-r2scan}).

\textbf{Correction.} The operator retrieves the pre-computed bias direction
$\hat b$ for the (element, family) pair from the distillation corpus, projects
the prediction onto the orthogonal complement of $\hat b$, and applies the
scalar correction $c = \langle y - y_{\mathrm{ref}}, \hat b \rangle / \|\hat b\|^2$.

\textbf{Output.} The corrected prediction $y_{\mathrm{corr}} = y - c \hat b$,
together with a conformal prediction interval \citep{shafer2008,angelopoulos2021}
constructed from the held-out empirical residuals of the calibration set.

\textbf{Transfer.} The correction is family-level --- it transfers to models
not yet trained, so long as they share the constraint. A model trained on a
new architecture but the same MatPES-PBE functional should receive the same
\texttt{mlip-pbe} correction vector; the 3$\times$3$\times$3 benchmark below
tests whether that transfer holds for the TensorNet-family QET checkpoint.

\textbf{Conformal uncertainty quantification.} Every corrected prediction carries
a nominal coverage guarantee when the conformal score is calibrated on the
held-out residuals of the same chemical class. The conformal score is the norm
of the residual in the bias-orthogonal subspace, so the uncertainty interval is
narrow when the model lands near the learned constraint surface and wide when it
departs from it. The Round-2 benchmark below reports the empirical coverage.

The claim survives the obvious circularity objection in the classical layer:
fitting the direction with the corrected model \emph{held out} removes a median
$69\%$ of squared elastic error out-of-sample (154 leave-one-model-out trials,
IQR $35$--$92\%$). The 3$\times$3$\times$3 MLIP benchmark below provides the
analogous Layer-2 test.

\section{Round-2 3$\times$3$\times$3 Benchmark and Operator Validation}
\label{sec:round2}

To test the operator at Layer~2 we ran a complete 3$\times$3$\times$3
elastic-constant benchmark for 16 cubic elemental metals (Ag, Al, Au, Ca, Cr,
Cu, Fe, Mo, Nb, Ni, Pd, Pt, Sr, Ta, V, W) and four MatPES 2025.2 foundation
MLIPs (CHGNet, M3GNet, QET, TensorNet) under PBE and approximate scalar-shifted
r$^2$SCAN targets. The matrix contains 128 cases and costs less than one CPU
core-hour.

\textbf{Raw accuracy.} Table~\ref{tab:raw} reports mean C$_{ij}$ MAE. QET is
the most accurate model on both functionals; CHGNet is the least. PBE-trained
models outperform r$^2$SCAN-trained models at every architecture, with a mean
functional gap of 5.7~GPa.

\begin{table}[h]
\centering
\caption{Raw mean C$_{ij}$ MAE (GPa) by model and functional in the 3$\times$3$\times$3 benchmark.}
\label{tab:raw}
\begin{tabular}{lrrr}
\toprule
Model & PBE & r$^2$SCAN & Overall \\
\midrule
CHGNet & 17.90 & 27.94 & 22.92 \\
M3GNet & 14.13 & 20.71 & 17.42 \\
TensorNet & 14.61 & 18.54 & 16.58 \\
QET & 13.41 & 15.46 & 14.44 \\
\textbf{All} & \textbf{15.01} & \textbf{20.66} & \textbf{17.84} \\
\bottomrule
\end{tabular}
\end{table}

\textbf{LOO correction operator.} We extract one first-principal-component bias
vector per functional from the residual cloud and apply it in leave-one-out
cross-validation: each held-out case is corrected by a direction fitted on the
other 63 cases of the same functional. The LOO operator reduces the overall
mean MAE from 17.84~GPa to 10.36~GPa and improves every model on both
functionals (Table~\ref{tab:loo}), with zero no-harm violations on the
held-out cases. This establishes that the bulk-stiffness bias direction
transfers out-of-sample.

\begin{table}[h]
\centering
\caption{Raw versus LOO-corrected mean C$_{ij}$ MAE (GPa).}
\label{tab:loo}
\begin{tabular}{lrrrrrr}
\toprule
Model & PBE raw & PBE corr. & r$^2$SCAN raw & r$^2$SCAN corr. & Overall raw & Overall corr. \\
\midrule
CHGNet & 17.90 & 11.01 & 27.94 & 13.57 & 22.92 & 12.29 \\
M3GNet & 14.13 & 8.37 & 20.71 & 11.82 & 17.42 & 10.09 \\
QET & 13.41 & 9.22 & 15.46 & 8.69 & 14.44 & 8.95 \\
TensorNet & 14.61 & 8.97 & 18.54 & 11.22 & 16.58 & 10.09 \\
\textbf{All} & \textbf{15.01} & \textbf{9.39} & \textbf{20.66} & \textbf{11.32} & \textbf{17.84} & \textbf{10.36} \\
\bottomrule
\end{tabular}
\end{table}

\textbf{Pre-registered hypothesis tests (H1--H4).} Table~\ref{tab:h} gives the
Round-2 outcomes. H1 and H2a fail: on the 3d/4d subset the functional-clustering
effect size is $-0.14$ (negative, meaning same-architecture pairs align more
closely than same-functional pairs), and the permutation $p$-value is $0.13$.
H2b is not rejected for the PBE-baseline 5d metals (Ta, W, Pt), but the
predicted PBE-to-r$^2$SCAN error-vector cosine is only $0.53$, below the
registered $0.8$ threshold. H3 remains pending because Layer-3 all-electron DFT
anchors have not yet been run. H4 is mixed: a single QET plus the LOO operator
beats the raw four-model ensemble on MAE (8.95 vs 12.62~GPa) and on exact MSE
(167.9 vs 288.7~GPa$^2$), but the 90\% conformal intervals achieve only
$86.7\%$ empirical coverage, below the registered $90\%$ threshold.

\begin{table}[h]
\centering
\caption{Round-2 pre-registered hypothesis outcomes.}
\label{tab:h}
\begin{tabular}{lp{8cm}l}
\toprule
Test & Prediction / kill condition & Status \\
\midrule
H1 & Effect size $\ge 0.30$ on 3d/4d; kill if $<0.20$ & \textbf{Kill triggered} ($-0.14$) \\
H2a & 3d/4d clusters by functional ($p<0.05$); kill if not & \textbf{Kill triggered} ($p=0.13$) \\
H2b & 5d (Ta,W,Pt) does not cluster ($p>0.20$) and cos$>0.8$ & Not killed; cos$=0.53$ misses threshold \\
H3 & Layer-2 XC bias aligns with Layer-3 DFT error ($\cos>0.5$) & Pending \\
H4 & Operator MSE $<$ ensemble MSE and coverage $\ge90\%$ & MAE/MSE pass; coverage fails \\
\bottomrule
\end{tabular}
\end{table}

\textbf{Interpretation.} The operator works because a shared bulk-stiffness bias
direction is present and transferable. The failure of the functional-clustering
predictions on this benchmark indicates that the effective binding constraint in
the 3$\times$3$\times$3 data is architecture family (or bonding class) rather
than training functional. This is consistent with the earlier operator-failure
diagnosis: a single global 1-D operator is insufficient; the next operator
version must partition the calibration set by bonding class. A class-aware LOO
operator lowers the oracle MAE further, from 10.36~GPa globally to 9.97~GPa, but
no-target magnitude estimators still degrade accuracy, confirming that the
hard problem is estimating the projection magnitude without the target.

\section{Meta-Science: Recursive Distillation of the Published MD Record}
\label{sec:metascience}

The projection law is not only a tool for new simulations; it is a lens for
reading the published molecular dynamics record. Across 559 classical
potentials, 8 foundation MLIPs, and 12 DFT implementations, the error
geometry is conserved and its direction rotates to the next constraint upstream
in the epistemic stack (Table~\ref{tab:layers}).

\begin{table}[t]
\centering
\caption{The projection law across one epistemic stack. Each row: an ensemble
of models, the constraint the law identifies as binding, the clustering
evidence (constraint vs.\ implementation), and the registered refutation
condition's status.}
\label{tab:layers}
\small
\begin{tabularx}{\textwidth}{@{}l l l X l@{}}
\toprule
Layer & Ensemble & Binding constraint & Evidence & Kill condition \\
\midrule
Classical IPs & 559 potentials & functional form & within-family $r{=}0.95$; PR invariant 40 yr & (observational) \\
Foundation MLIPs & $4{\times}2$ MatPES $+$ anchors & training functional & $S_{\mathrm{func}}{=}{+}0.317$ vs ${-}0.093$, exact perm.\ $p{=}0.029$ & not triggered \\
DFT codes & 12 ACWF methods & pseudopotential table & $S_{\mathrm{table}}{=}{+}0.526$ vs ${+}0.265$, $p{=}0.017$ & not triggered \\
\bottomrule
\end{tabularx}
\end{table}

Between layers, the direction \emph{rotates}: no transfer of classical
cross-family alignment to MLIP alignment is detectable (Spearman $\rho = 0.26$,
$p = 0.34$ across elements --- a low-powered test, so this is absence of
evidence for transfer rather than evidence of independence; the rotation claim
rests on the factorial clustering, not on this null), because the constraint
moved from functional form to training functional; and the MLIP-layer bias
direction is precisely the reference-theory error that layer 3 isolates. Within
layers, constraints \emph{nest}: when the registered constraint is removed
(r$^2$SCAN for PBE; plane-wave pseudopotential sharing for SIESTA), residual
alignment reveals the next constraint in line (5d correlation physics; basis
set). Anisotropy is conserved throughout, in the operational sense that the
participation ratio stays far below the ambient dimension at every layer ---
and the quantitative agreement across layers is striking: median PR $= 1.09$
(classical, 42 multi-element potentials), $1.13$--$2.09$ per element
(foundation MLIPs, $n = 8$--$11$ models), and median $1.10$ (range $[1.00, 2.11]$,
rank-one share median $0.95$) across 134 above-floor ACWF systems with
$\geq 6$ methods. At no layer, in no ensemble, under no intervention did errors
approach isotropy, as Theorem~\ref{thm:gauge} requires of any family whose
systematic error dominates its fitting scatter.

The closest antecedents are the persistence of shared biases across climate
model generations \citep{tian2020,knutti2013}, the convergence of training
trajectories onto shared low-dimensional manifolds \citep{mao2024}, and
representation convergence across architectures \citep{huh2024}; none states
conservation-plus-rotation as a law, and none possesses the factorial
constraint-vs-implementation evidence presented here.

\textbf{Implications.} \textbf{Calibration without retraining.} Because the
shared error is one direction, one correction vector per (element, observable)
repairs every model in the family at once --- and the claim survives the obvious
circularity objection: fitting the direction with the corrected model
\emph{held out} removes a median $69\%$ of squared elastic error out-of-sample.

\textbf{A trust rule for ensembles.} Cross-model alignment is a zero-cost
diagnostic: high alignment $\Rightarrow$ the ensemble is constraint-bound,
calibratable, and its spread \emph{understates} uncertainty; low alignment at
the noise floor $\Rightarrow$ genuine convergence (the Ni control); low
alignment far above the floor $\Rightarrow$ the constraint is heterogeneous
(magnetic BCC metals; difficult elements across pseudopotential tables) and no
shared correction exists.

\textbf{Benchmark design.} Leaderboards that pool across families average away
the very direction that identifies what to fix. Reporting error
\emph{directions} stratified by candidate constraints --- the factorial
pattern used here --- converts benchmarking from scoring into diagnosis.

\textbf{Reading consensus.} Where ground truth is delayed (climate
projections, materials discovery), inter-model agreement should be priced as a
measurement of shared constraint; the factorial designs show how to estimate
\emph{which} constraint, today, from models alone.

\section{Limitations}
\label{sec:limits}

The formal chain covers convex reachable sets and deterministic
bias-plus-noise spectra; smooth non-convex families (local normal cones) and
the finite-sample step from noisy ensembles to the exact second moment are
standard but unformalized. The MLIP layer uses a stress/strain elasticity harness on sixteen cubic
elemental metals and elastic observables --- the lowest-dimensional, most
symmetric corner of materials space --- with 0~K published DFT-PBE references
and approximate scalar-shifted r$^2$SCAN targets. The measured residual
therefore mixes fitting error, exchange--correlation bias, reference-standard
differences, and the r$^2$SCAN scalar-shift approximation. For the magnetic
elements (Fe, Cr, Ni, V) the spin protocol is a live confound. The Round-2
pre-registered functional-clustering tests (H1, H2a) did not survive on the
3d/4d subset; this is reported as a failure, not suppressed. The operator
results are an oracle ceiling or a class-aware upper bound; a fully deployable
no-target operator remains to be validated.
The DFT layer inherits ACWF's protocol choices, and the SIESTA nested-constraint
reading is post-hoc until an independent localized-basis code is tested.
Statistically: the permutation lattices are small (minimum attainable $p = 1/70$),
seven registered predictions across two experiments invite multiplicity concerns
that round 2's single-primary-endpoint design addresses; in the present
manuscript the two kill conditions were primary (cluster-by-architecture for
MLIPs, cluster-by-code for DFT), while the other five predictions were
auxiliary robustness checks with no formal hierarchical correction applied, and
the registered kill conditions were asymmetric (refutation required a sign
reversal; round 2 registers a symmetric equivalence bound). Four of the seven
registered predictions failed and are reported as failures (2/4 and 1/3
passing per experiment); the nested-constraint attributions of those misses
are hypotheses for the next registration round, not confirmed claims.

\section{Reproducibility}
\label{sec:repro}

Both experiments were pre-registered with thresholds and refutation
conditions committed to version control before model execution for the
$4{\times}2$ (commit \texttt{dffbe595}) and before any error-vector statistic
was computed for the ACWF analysis (commit \texttt{ebf39e33}); all factorial
statistics are replayable from a two-tier kit in which Tier~1 verifies the
factorial-experiment statistics from committed raw data using only
NumPy/SciPy in seconds (classical-layer and robustness statistics are
reproduced by accompanying scripts in the same kit), and Tier~2 re-derives
the raw elastic constants from public checkpoints through a frozen harness
(deterministic within 1~GPa; the Gate-0 harness regression is bit-exact).
The complete kit, pinned datasets, and pre-registrations are publicly served at
\url{https://storage.googleapis.com/shed-489901-replication/error-geometry/v1-10c18ace/index.html};
the citable Zenodo snapshot is archived at
\url{https://doi.org/10.5281/zenodo.20787874}. The Lean~4 artifact (225 theorem
and lemma declarations in the current build, including the seven theorems of
this paper; zero \texttt{sorry}, zero new axioms, pinned Mathlib) and the
theorem-to-statistic contract accompany the kit. The harness was validated by
independent energy- vs.\ stress-method agreement ($\le 2.4\%$), bit-exact
regression, and strain-window stability. ACWF data are from the public
archive of \citet{bosoni2024}; MatPES checkpoints from \citet{kaplan2025}.

\section*{Acknowledgments}
This work was supported by Lupine Science. AI assistance (Anthropic Claude,
OpenAI Codex, and Kimi) was used under the author's direction for statistical
verification and recomputation, Lean proof engineering, literature triage, and
drafting; all scientific claims and decisions were reviewed and approved by the
author, who takes sole responsibility for the content.

\bibliographystyle{plainnat}
\bibliography{references}

\end{document}
